\section{A quantitative geometric extraction on every finite itinerary}
\label{sec:geometric-raw-extraction}
The power--logarithm extraction gives a residual at each finite packet,
but does not bound its derivative norm as the collision cutoff grows.
We construct a different exact extraction. Its removed measure has a
quantified total variation, and its residual has a uniform, explicit
second-derivative budget. The construction takes place before coarea,
on the physical collision cylinder. In particular it neither estimates
individual prepared jets nor assumes separated critical values.

Write $L_m=S_{m,R}\tau_R$, and use $(\alpha,p)$ throughout for derivatives
and area. All constants in this section are uniform in $R\in I$.
Put
\[
 R_-=9/20,\qquad R_+=47/100,\qquad \ell_0=3/50,
 \qquad H_0=1+R_+/\ell_0=53/6.
\]
The integer $m$ is a number of physical collisions, not returns.

\subsection{A direction transverse to every regular roof level}
Lemma~\ref{lem:self-contained-collision-differential} gives a
self-contained derivation of the coordinate differential and action
used below, with both signs and the determinant in invariant coordinates.
\begin{lemma}[A uniform nonstationary direction]\label{lem:roof-transverse-direction}
On a regular $m$-collision branch, $m\ge1$, write
\[
 D T_R^m=\begin{pmatrix}A_m&B_m\\ C_m&D_m\end{pmatrix}.
\]
Then $B_m\ne0$ and $0<A_m/B_m\le H_0$. The total roof satisfies
\begin{equation}\label{eq:uniform-roof-gradient}
 |\nabla L_m|\ge \frac{R_-}{\sqrt{1+H_0^2}}|p|.
\end{equation}
On $p\ne0$ the vector field
\begin{equation}\label{eq:roof-transverse-field}
 X_m=-\frac1{Rp}\left(\partial_\alpha-
                      \frac{A_m}{B_m}\partial_p\right)
 \quad\hbox{satisfies}\quad X_m L_m=1.
\end{equation}
No incidence bound at the interior collisions is required for these
three assertions.
\end{lemma}
\begin{proof}
Let $c=\sqrt{1-p^2}$, $c'=\sqrt{1-(p')^2}$ and $v=\tau_R/R$ for
one regular flight. In the coordinates used here its differential is
\begin{equation}\label{eq:positive-collision-matrix}
 D T_R=-\begin{pmatrix}
 (v+c)/c'&v/(cc')\\
 v+c+c'&(v+c')/c
 \end{pmatrix}=:-P.
\end{equation}
One obtains this formula by differentiating the straight flight and
its reflection: in arc length and outgoing angle the four entries,
before multiplication by $-1/c'$, are
$\tau_R/R+c$, $\tau_R$, $\tau_R/R^2+(c+c')/R$, and
$\tau_R/R+c'$. Change variables by $r=R\alpha$, $p=\sin\varphi$.
This is also the fixed-table specialization of the differential in
\cite[Section 3.1]{SYZ}. Its determinant in $(\alpha,p)$ is one.

All entries of $P$ are positive. The ratios of the first to the second
entry in either row are
\[
 c+\frac{c^2}{v},\qquad c+\frac{c^2}{v+c'},
\]
respectively, and are at most $H_0$. The first row of the positive
product $P_{m-1}\cdots P_0$ is a positive linear combination of the
two rows of $P_0$. Its entry ratio is consequently a weighted average
of these two ratios, with positive weights. This proves the assertion
about $A_m/B_m$, including $m=1$.

The exact action identity of Lemma~\ref{lem:exact-collision-action}
gives
\[
 \partial_\alpha L_m=R(p_m A_m-p),\qquad
 \partial_p L_m=Rp_m B_m.
\]
Subtracting $A_m/B_m$ times the second identity from the first gives
$-Rp$. Cauchy--Schwarz proves \eqref{eq:uniform-roof-gradient}; the
same calculation proves \eqref{eq:roof-transverse-field}.
\end{proof}

\subsection{A one-collision regularity cutoff}
\begin{lemma}[Polynomial sublevel sets in a rectangle]
\label{lem:elementary-polynomial-sublevel}
Fix $D\ge1$ and $c_0>0$. A real polynomial $P$ in two variables, of
degree at most $D$ in each variable and with maximum coefficient
modulus at least $c_0$, satisfies
\begin{equation}\label{eq:elementary-polynomial-sublevel}
 |\{(x,y)\in[-1,1]^2:|P(x,y)|\le h\}|
                  \le C_{D,c_0}h^{1/(2D)},\qquad 0<h<1.
\end{equation}
An upper bound on the other coefficients is not needed.
\end{lemma}
\begin{proof}
For a one-variable polynomial of degree at most $D$, a sublevel set of
length $s>0$ contains $D+1$ points whose pairwise distances are at least
$c_Ds$. For example choose successive quantiles of its restricted
Lebesgue measure, discarding a small part at the endpoints. Lagrange
interpolation on these points bounds every coefficient by $C_D h s^{-D}$
when the polynomial is at most $h$ there. It follows that a coefficient
of modulus at least $b>0$ gives the one-variable bound
$C_D(h/b)^{1/D}$; when this exceeds two the bound is automatic.

Write $P(x,y)=\sum_j a_j(y)x^j$ and choose $a_j$ with a coefficient
of modulus at least $c_0$. The set where $|a_j(y)|\le h^{1/2}$ has
length at most $C_{D,c_0}h^{1/(2D)}$. On its complement, apply the
one-variable bound to $P(\cdot,y)$ with $b=h^{1/2}$. Integrating over
$y$ proves \eqref{eq:elementary-polynomial-sublevel}. Constant
polynomials cause no difficulty in either step.
\end{proof}

Fix an even $\theta\in C^\infty(\R;[0,1])$, zero on $[-1,1]$ and one
outside $(-2,2)$. A finite, parameter-independent set $\mathcal D$ of
nonzero lattice centers contains every possible next collision center,
by the uniform finite horizon. For $d\in\mathcal D$ put
\begin{equation}\label{eq:all-candidate-discriminants}
 \Delta_{d,R}(\alpha,p)
   =R^2-\bigl((d-Rn_\alpha)\mathbin{\times}
          (\sqrt{1-p^2}n_\alpha+pt_\alpha)\bigr)^2.
\end{equation}
The cross denotes the scalar planar cross product. We use all candidate
discriminants, including candidates that are not selected by the actual
flight. This prevents a discontinuity from being missed on the side
where a tangent disk ceases to be the first collision.
Lemma~\ref{lem:complete-first-hit-guard} specifies the finite candidate
set, four bounded charts and uniform coefficient minimum, and proves
that every label transition is covered.

\begin{lemma}[Uniform regularity cutoffs]\label{lem:uniform-collision-cutoff}
There are $\varepsilon_0>0$ and smooth functions
$\chi_{R,\varepsilon}:M\to[0,1]$, $0<\varepsilon\le\varepsilon_0$,
vanishing on neighborhoods of the one-collision singularities, grazing
endpoints and every boundary of $Y_R^*$, such that, with $\kappa_0=1/16$,
\begin{equation}\label{eq:collision-cutoff-measure}
 \nu\{\chi_{R,\varepsilon}\ne1\}\le C\varepsilon^{\kappa_0}.
\end{equation}
On the support of this cutoff and its first three derivatives, every
selected one-collision branch and roof, as well as the cutoff itself,
have $C^3(\alpha,p)$ bounds at most $(C/\varepsilon)^C$.
The actual section membership and lattice label are locally constant
where needed. Values at the excluded singularities are fixed to zero.
\end{lemma}
\begin{proof}
Let $\mathcal A_R$ contain the angular endpoints and $\mathcal B_R$ the
$p$-endpoints of all the finitely many rectangles defining $Y_R^*$.
Set $c(p)=\sqrt{1-p^2}$ and take the product
\begin{equation}\label{eq:explicit-collision-cutoff}
 \theta(c(p)/\varepsilon)
 \prod_{d\in\mathcal D}\theta(\Delta_{d,R}/\varepsilon)
 \prod_{a\in\mathcal A_R}\theta(\sin(\alpha-a)/\varepsilon)
 \prod_{b\in\mathcal B_R}\theta((p-b)/\varepsilon).
\end{equation}
The extra angular zeros at antipodal points only enlarge the removed
set by finitely many strips. All factors are globally defined on the
angular circle. The first factor makes the product identically zero
near $p=\pm1$, so its zero extension is smooth at these endpoints.

Here is an explicit uniform sublevel estimate for each discriminant.
On a bounded rational chart for $\alpha$, write $t$ for its half-angle
coordinate, and put $s=\tan(\varphi/2)$, so that $|s|\le1$,
$c=(1-s^2)/(1+s^2)$ and $p=2s/(1+s^2)$. The cross product in
\eqref{eq:all-candidate-discriminants} has denominator
$(1+t^2)(1+s^2)$ and numerator of total degree at most four. Thus
$\Delta_{d,R}$ has a bounded positive denominator and a numerator of
total degree at most eight. That numerator is never the zero
polynomial: at any fixed initial boundary point, $|d-Rn_\alpha|>R$,
and its cross product with a varying unit direction cannot have
constant square $R^2$ on an interval of directions. Compactness of $I$
and finiteness of $\mathcal D$ and the charts give a positive lower
bound on the maximum coefficient modulus. Apply
Lemma~\ref{lem:elementary-polynomial-sublevel} with $D=8$.
The coordinate Jacobian of $\nu$ in these bounded $(t,s)$ charts is
bounded above. This proves the bound $C\varepsilon^{1/16}$ for the
discriminant factors. The remaining factors remove finitely many
strips of measure $O(\varepsilon)$, or $O(\varepsilon^2)$ at grazing.
A union bound proves \eqref{eq:collision-cutoff-measure}.

Every possible change of first-hit label occurs at a candidate
tangency: a positive entry root cannot cross zero because distinct
disks are separated, and two distinct disks cannot have the same
entry point at the same positive time. Those tangencies are among the
zero sets in \eqref{eq:all-candidate-discriminants}. On a selected
branch with $\chi_{R,\varepsilon}>0$, its entry root is
\[
 (d-Rn_\alpha)\cdot v-\sqrt{\Delta_{d,R}},
 \qquad \Delta_{d,R}>\varepsilon,
\]
with $c(p)>\varepsilon$. Differentiating this formula and the Euclidean
reflection formula through order three gives a fixed polynomial loss
in $\varepsilon^{-1}$. The outgoing point lies on a circle of radius
at least $R_-$, so local angular coordinates add only fixed losses.
Changing from $\varphi$ to $p$ costs further fixed powers of
$c(p)^{-1}$. Differentiating the factors of
\eqref{eq:explicit-collision-cutoff} has the same property. The estimate
also holds on their derivative supports, because a smooth factor which
vanishes on an interval is flat at its endpoints. Section membership
can change only on the explicitly removed rectangle boundaries.
\end{proof}

For $L\ge1$ define, on regular orbits and then by zero extension,
\begin{equation}\label{eq:whole-itinerary-cutoff}
 \Gamma_{L,R,\varepsilon}(x)
  =\theta(p(x)/\varepsilon)
                    \prod_{j=0}^{L}\chi_{R,\varepsilon}(T_R^jx).
\end{equation}
The first factor removes a thin normal strip at the initial state.
This one strip, rather than neighborhoods chosen separately at every
critical value, suffices for \eqref{eq:uniform-roof-gradient}.

\begin{lemma}[Whole-itinerary bounds]\label{lem:whole-itinerary-cutoff}
The function in \eqref{eq:whole-itinerary-cutoff} is globally $C^3$
after zero extension across its excluded boundaries, and
\begin{align}
 \int(1-\Gamma_{L,R,\varepsilon})\dd\nu_R^*
          &\le C(L+1)\varepsilon^{1/16},
             \label{eq:whole-itinerary-mass}\\
 \|\Gamma_{L,R,\varepsilon}\|_{C^3}
          &\le \exp\{C(L+1)\log(C/\varepsilon)\}.
             \label{eq:whole-itinerary-derivatives}
\end{align}
On its derivative supports, for $1\le m\le L$, the field $X_m$ has
$C^2$ norm bounded by the right side of
\eqref{eq:whole-itinerary-derivatives}, after changing $C$.
\end{lemma}
\begin{proof}
Use $1-\prod b_j\le\sum(1-b_j)$ for $0\le b_j\le1$, invariance of
$\nu$, and $\nu_R^*\le c_*^{-1}\nu$. The initial normal strip has
measure $O(\varepsilon)$, and each of the other $L+1$ terms is bounded
by \eqref{eq:collision-cutoff-measure}. This proves
\eqref{eq:whole-itinerary-mass} without counting itinerary cells.

On a nonzero derivative support all earlier branches are regular and
all the one-step bounds of Lemma~\ref{lem:uniform-collision-cutoff}
apply. The chain rule through fixed order three, iterated $j$ times,
gives at most $(C/\varepsilon)^{C(j+1)}$. The product rule for $L+2$
factors contributes at most a fixed power of $L+2$, absorbed by
$\exp(C(L+1))$. At any boundary of a finite itinerary, an earlier
cutoff factor vanishes on a neighborhood on both sides. At a section
boundary the corresponding factor at that time does the same.
Consequently these piecewise derivatives extend by zero and give the
global estimate.

Every entry of the positive matrix in
\eqref{eq:positive-collision-matrix} is at least
$v_0=\ell_0/R_+=6/47$. In particular $|B_m|\ge v_0^m$. The quotient
rule, the $C^3$ bounds for $T_R^m$, and this lower bound give a
$C^2$ bound for $A_m/B_m$ of the same exponential form. On the relevant
supports $|p|\ge\varepsilon$. Formula~\eqref{eq:roof-transverse-field}
therefore gives the assertion about $X_m$.
\end{proof}

\subsection{An exact residual for bounded-variation return marks}
Fix $q\ge0$, indices $0\le k_1,\ldots,k_q\le n$, and complex section
functions $a_1,\ldots,a_q$ with bounded supremum and BV zero extensions
$a_i^0$ to $M$. The number $q$ is fixed; repeated marks are allowed.
Put
\begin{equation}\label{eq:geometric-weight-budgets}
 \begin{split}
 w(x)&=\prod_{i=1}^q a_i((F_R^*)^{k_i}x),\\
 M&=\prod_{i=1}^q\|a_i\|_\infty,\qquad
 V=\sum_{i=1}^q\|a_i^0\|_{\mathrm{BV}}
                         \prod_{j\ne i}\|a_j\|_\infty.
 \end{split}
\end{equation}
For $q=0$ these conventions give $w=1$, $M=1$ and $V=0$.
All measures below use the normalized section law $\nu_R^*$ and retain
the exact weight $w$. For the earlier unnormalized single-mark convention,
choose $w=c_*a\circ(F_R^*)^k$.

\begin{theorem}[All-itinerary geometric extraction]
\label{thm:geometric-extraction-budget}
For $1\le n\le L$ and $0<\varepsilon\le\varepsilon_0$, there is an exact
decomposition of the actual finite-count measure
\begin{equation}\label{eq:geometric-exact-decomposition}
 \mu_{n,R}^{w,L}=E_{n,R,\varepsilon}^{w,L}
                         +Q_{n,R,\varepsilon}^{w,L}
\end{equation}
with compactly supported mixed densities $e_\varepsilon$ and
$q_\varepsilon$, respectively. The residual satisfies
\begin{align}
 \|Q_{n,R,\varepsilon}^{w,L}\|_{\TV}&\le M,\notag\\
 \|E_{n,R,\varepsilon}^{w,L}\|_{\TV}
     &\le C_q\{\varepsilon V+M(L+1)\varepsilon^{1/16}\},
                   \label{eq:geometric-extraction-mass}\\
 \mathcal A_{2,\varepsilon}(n,L,R,w)
  :=\sum_{\ell\in\Z^3}\|\partial_t^2q_\varepsilon(\ell,\cdot)\|_1
     &\le C_q M\exp\{C_q(L+1)\log(C_q/\varepsilon)\}.
                   \label{eq:geometric-derivative-budget}
\end{align}
Each residual component belongs to $W^{2,1}(\R)$. The constants do not
depend on $n$, the marks or the radius. When $w=1$, both terms of
\eqref{eq:geometric-exact-decomposition} are positive measures.
No subanalytic assumption on the marked functions is imposed.
\end{theorem}
\begin{proof}
Set $a_{i,\varepsilon}=\mathsf S_{\varepsilon/8}a_i^0$, using the
mean-preserving smoothing of Lemma~\ref{lem:bv-smoothing}, and let
$w_\varepsilon$ be the corresponding product at the same actual marks.
Differentiating the same smoothing kernel three times gives
$\|a_{i,\varepsilon}\|_{C^3}\le C\varepsilon^{-3}\|a_i\|_\infty$;
composition with $p=\sin\varphi$ is not needed for this flat-coordinate bound.
Define
\begin{equation}\label{eq:geometric-residual-definition}
 Q_{n,R,\varepsilon}^{w,L}=(J_{n,R})_*
 \bigl(\one_{\{N_{n,R}\le L\}}\Gamma_{L,R,\varepsilon}
                                     w_\varepsilon\nu_R^*\bigr),
 \qquad E_{n,R,\varepsilon}^{w,L}=\mu_{n,R}^{w,L}-Q_{n,R,\varepsilon}^{w,L}.
\end{equation}
An empty product is not smoothed. Supremum contraction gives
$|w_\varepsilon|\le M$. Telescoping the products and invariance of
$F_R^*$ give
\[
 \int|w-w_\varepsilon|\dd\nu_R^*
 \le c_*^{-1}\sum_i\|a_i^0-a_{i,\varepsilon}\|_{L^1(\nu)}
                        \prod_{j\ne i}\|a_j\|_\infty
 \le C_q\varepsilon V.
\]
Combining this with \eqref{eq:whole-itinerary-mass} proves the two
variation bounds. When $w=1$ the residual source weight is exactly
$\Gamma_{L,R,\varepsilon}\one_{\{N_n\le L\}}$, which lies between zero
and the source weight of $\mu^{1,L}_{n,R}$. This proves positivity.

It remains to establish the derivative estimate including the sum over
all components. Partition the initial section by its finite regular
collision itinerary through $L$, all its section decisions, and the
value $m=N_{n,R}\le L$. There are finitely many pieces, but no estimate
of their number will be used. On one such piece $D$, every mark is
$T_R^{j_i}x$ for a fixed $j_i\le m$, the lattice label
$\ell=(\sum_{j<m}\kappa_R\circ T_R^j,m)$ is constant, and the time
coordinate is $L_m(x)$. The density on the initial cylinder is
\[
 b_D=(4\pi c_*)^{-1}\one_D\Gamma_{L,R,\varepsilon}w_\varepsilon.
\]
Lemma~\ref{lem:source-partition-zero-extension} supplies the formal
partition and the earliest-forbidden-time argument. Its zero extension
across the boundaries of $D$ is $C^2$: all changes
of itinerary or section decision have already been removed on both
sides by \eqref{eq:whole-itinerary-cutoff}. Fixed-order chain and product
rules, smoothing in $C^3$, and Lemma~\ref{lem:whole-itinerary-cutoff}
give on $D$
\begin{equation}\label{eq:geometric-divergence-bound}
 |b_D|+|\operatorname{div}(b_DX_m)|+
 |\operatorname{div}(X_m\operatorname{div}(b_DX_m))|
       \le C_qM\exp\{C_q(L+1)\log(C_q/\varepsilon)\}.
\end{equation}
All vector-field products are extended by zero outside their regular
pieces; the orbit cutoffs vanish on neighborhoods of every excluded
boundary. A further smooth cutoff in $p$ can remove the factor $1/p$
away from the supports. No extension of the billiard map through a
singularity is used.
Angular integration is periodic, and all source functions vanish near
$p=\pm1$. Hence there are no coordinate boundary terms.

For a test function $\phi\in C_c^\infty(\R)$,
$X_m(\phi\circ L_m)=\phi'\circ L_m$. Integration by parts therefore
identifies the first and second distributional derivatives of the
pushforward density as the pushforwards of
\begin{equation}\label{eq:geometric-density-derivatives}
 \operatorname{div}(b_DX_m),\qquad
 \operatorname{div}(X_m\operatorname{div}(b_DX_m)),
\end{equation}
respectively. The signs are positive in both distributional identities.
On these supports $L_m$ is a submersion by
\eqref{eq:uniform-roof-gradient}, so coarea gives $L^1$ densities for
all three pushforwards. Their norms are at most the integrals of the
absolute source functions.

The interiors of the pieces $D$ are disjoint. Summing these source
integrals, even after summing over all output labels, pays only the
fixed area of the collision cylinder. It does not multiply
\eqref{eq:geometric-divergence-bound} by an itinerary count. This proves
\eqref{eq:geometric-derivative-budget} and the $W^{2,1}$ assertion.
The actual finite-count measure has a density for bounded measurable
weights by coarea on its regular branches; alternatively
\eqref{eq:uniform-roof-gradient} excludes a zero gradient outside the
null initial line $p=0$. Subtraction gives the density of $E$.
All time supports lie in $[0,L\tau_{\max}]$, and all lattice supports
are finite.
\end{proof}

\begin{remark}[What is extracted]\label{rem:geometric-extraction-scope}
The measure $E$ includes the discarded parts of every regular critical
word, every grazing and first-hit transition, every section-decision
boundary, and all their iterated preimages. Any singular output image
created by those parts remains in $E$. There is no unlisted branch in
\eqref{eq:geometric-exact-decomposition}. This extraction is not the
power--logarithm jet extraction of Section~\ref{sec:finite-count-extraction};
its new derivative bound does not bound the old jet residual. In
particular, small total variation of $E$ does not imply a small
essential supremum of $e_\varepsilon$.
\end{remark}
